% =============================================================================
% Chapter 8: Mathematical Foundations
% =============================================================================
%
% Version: 1.0 (January 2026)
%
% Purpose: Formal framework
%
% Primary Appendix: See appendix references below
% Secondary Appendices: G (Glossary), F (Examples)
%
% Prerequisites: Chapters 1-7
% Leads to: Chapter 9
%
% Chapter Type: B (Foundation)
% =============================================================================


%%%%%%%%%%%%%%%%%%%%%%%%%%%%%%%%%%%%%%%%%%%%%%%%%%%%%%%%%%%%%%%%%%%%%%%%%%%%%%%
% QUICK REFERENCE BOX
%%%%%%%%%%%%%%%%%%%%%%%%%%%%%%%%%%%%%%%%%%%%%%%%%%%%%%%%%%%%%%%%%%%%%%%%%%%%%%%
\begin{tcolorbox}[
    colback=gray!5!white,
    colframe=gray!75!black,
    title={\textbf{Begriffe in diesem Kapitel}},
    fonttitle=\bfseries
]
\small
\begin{tabular}{@{}ll@{}}
\textbf{Key Concept} & \textbf{Definition} \\
\midrule
Formal framework & See Section \ref{sec:ch8-intro} \\
\end{tabular}
\end{tcolorbox}



%%%%%%%%%%%%%%%%%%%%%%%%%%%%%%%%%%%%%%%%%%%%%%%%%%%%%%%%%%%%%%%%%%%%%%%%%%%%%%%
% INTUITION BOX
%%%%%%%%%%%%%%%%%%%%%%%%%%%%%%%%%%%%%%%%%%%%%%%%%%%%%%%%%%%%%%%%%%%%%%%%%%%%%%%
\begin{tcolorbox}[
    colback=blue!5!white,
    colframe=blue!75!black,
    title={\textbf{Intuition: A Motivating Example}},
    fonttitle=\bfseries
]
Consider \textbf{Anna}, a professional facing a decision that involves formal framework.

She initially evaluates the choice using standard criteria, but something feels incomplete.
\textbf{Thomas}, her colleague, suggests considering additional dimensions beyond the obvious.

This chapter explores what Anna discovers when she applies the EBF framework---how
formal framework interact with broader welfare considerations.

\medskip
\textit{By the end of this chapter, you will understand why Anna's initial analysis
was incomplete and how the EBF approach provides a more comprehensive view.}
\end{tcolorbox}



%%%%%%%%%%%%%%%%%%%%%%%%%%%%%%%%%%%%%%%%%%%%%%%%%%%%%%%%%%%%%%%%%%%%%%%%%%%%%%%
% CENTRAL QUESTION BOX
%%%%%%%%%%%%%%%%%%%%%%%%%%%%%%%%%%%%%%%%%%%%%%%%%%%%%%%%%%%%%%%%%%%%%%%%%%%%%%%
\begin{tcolorbox}[
    colback=red!5!white,
    colframe=red!75!black,
    title={\textbf{Central Question}},
    fonttitle=\bfseries
]
\Large
\centering
\textit{How does formal framework contribute to understanding human welfare and decision-making?}
\end{tcolorbox}



%%%%%%%%%%%%%%%%%%%%%%%%%%%%%%%%%%%%%%%%%%%%%%%%%%%%%%%%%%%%%%%%%%%%%%%%%%%%%%%
% CHAPTER OVERVIEW
%%%%%%%%%%%%%%%%%%%%%%%%%%%%%%%%%%%%%%%%%%%%%%%%%%%%%%%%%%%%%%%%%%%%%%%%%%%%%%%
\subsection*{Chapter Overview}

This chapter is organized as follows:

\begin{itemize}
\item \textbf{Section \ref{sec:ch8-intro}}: Introduction and motivation
\item \textbf{Section \ref{sec:ch8-theory}}: Theoretical foundations
\item \textbf{Section \ref{sec:ch8-applications}}: Applications and examples
\item \textbf{Section \ref{sec:ch8-summary}}: Summary and conclusions
\end{itemize}




%%%%%%%%%%%%%%%%%%%%%%%%%%%%%%%%%%%%%%%%%%%%%%%%%%%%%%%%%%%%%%%%%%%%%%%%%%%%%%%
% CHAPTER SCOPE BOX
%%%%%%%%%%%%%%%%%%%%%%%%%%%%%%%%%%%%%%%%%%%%%%%%%%%%%%%%%%%%%%%%%%%%%%%%%%%%%%%
\begin{tcolorbox}[
    colback=purple!5!white,
    colframe=purple!75!black,
    title={\textbf{Chapter Scope: Ziel / In-Scope / Out-of-Scope}},
    fonttitle=\bfseries
]

\textbf{Ziel (Objective):}
Develop the conceptual understanding of formal framework as a foundation for the EBF framework.

\vspace{0.3cm}
\textbf{In-Scope:}
\begin{itemize}[nosep]
    \item Conceptual introduction to formal framework
    \item Motivation and intuition
    \item Connection to subsequent chapters
    \item Key definitions and concepts
\end{itemize}

\vspace{0.3cm}
\textbf{Out-of-Scope (delegiert an Appendices):}
\begin{itemize}[nosep]
    \item Complete formal treatment $\rightarrow$ FORMAL appendices
    \item Empirical validation $\rightarrow$ METHOD appendices
    \item Domain applications $\rightarrow$ Application chapters
\end{itemize}

\end{tcolorbox}

%%%%%%%%%%%%%%%%%%%%%%%%%%%%%%%%%%%%%%%%%%%%%%%%%%%%%%%%%%%%%%%%%%%%%%%%%%%%%%%
% APPENDIX REFERENCES BOX
%%%%%%%%%%%%%%%%%%%%%%%%%%%%%%%%%%%%%%%%%%%%%%%%%%%%%%%%%%%%%%%%%%%%%%%%%%%%%%%
\begin{tcolorbox}[
    colback=yellow!10!white,
    colframe=orange!75!black,
    title={\textbf{Appendix References for This Chapter}},
    fonttitle=\bfseries
]
\small
\begin{tabular}{@{}lp{8cm}@{}}
\textbf{Appendix} & \textbf{Content} \\
\midrule
D (FORMAL-CSTAR) & Complete mathematical proofs for $C^*$ \\
A (FORMAL-LIMITS) & Derivation of limiting cases \\
G (REF-GLOSSARY) & Symbol definitions and terminology \\
\end{tabular}

\medskip
\textit{For complete proofs and derivations, see the referenced appendices.}
\end{tcolorbox}

\section{Mathematical Formalization}
\label{sec:mathformalization}
%%%%%%%%%%%%%%%%%%%%%%%%%%%%%%%%%%%%%%%%%%%%%%%%%%%%%%%%%%%%%%%%%%%%%%%%%%%%%%%

The preceding sections developed intuition through examples.
This section provides the mathematical machinery (for standard microeconomic foundations, see \citealt{mascolell1995microeconomic}) that makes the framework precise,
testable, and applicable. We build up from simple two-dimensional models
to the full 144-component structure, always keeping the intuition visible.

\subsection{The Core Insight: Why Two Dimensions?}

\paragraph{The Roberts Reduction.}
\citet{roberts2004} showed that organizational choices---however numerous---
cluster along two fundamental dimensions:
\begin{itemize}
  \item \textbf{What we want} (goals, strategy, aspirations)
  \item \textbf{How we organize} (structure, processes, constraints)
\end{itemize}

This insight generalizes beyond organizations.
At every level, agents face a version of the same choice:
\begin{center}
\renewcommand{\arraystretch}{1.3}
\begin{tabular}{lll}
\hline
\textbf{Level} & \textbf{``What we want''} & \textbf{``How we organize''} \\
\hline
Individual & Life goals, values & Lifestyle, habits, routines \\
Household & Family priorities & Role allocation, rules \\
Organization & Business strategy & Structure, processes \\
Region & Economic positioning & Institutions, infrastructure \\
Nation & Societal model & State organization, laws \\
Global & World order vision & Governance architecture \\
\hline
\end{tabular}
\end{center}

\paragraph{The Two Poles.}
Each dimension spans between two poles:
\begin{itemize}
  \item \textbf{Strategy axis:} Stability/security ($s = 0$) vs. Dynamism/growth ($s = 1$)
  \item \textbf{Structure axis:} Hierarchy/control ($\sigma = 0$) vs. Flexibility/autonomy ($\sigma = 1$)
\end{itemize}

\paragraph{Example: National Models.}
\begin{itemize}
  \item Germany ($s \approx 0.3$, $\sigma \approx 0.3$): Stability-oriented, rule-based, coordinated
  \item USA ($s \approx 0.7$, $\sigma \approx 0.7$): Growth-oriented, flexible, market-based
  \item Both coherent ($s \approx \sigma$), different configurations
\end{itemize}

\subsection{The Two-Axis Model: Formal Definition}

\begin{definition}[Configuration Space]
For each level $\ell \in \mathcal{L} = \{Ind, HH, Org, Reg, Nat, Glob\}$,
the configuration space is the unit square $\mathcal{S} = [0,1] \times [0,1]$
with coordinates $(s, \sigma)$ representing strategy and structure.
\end{definition}

\begin{definition}[Strategy and Structure Axes]
\begin{align}
  s &\in [0,1] \quad \text{(Strategy: stability $\leftrightarrow$ dynamism)} \\
  \sigma &\in [0,1] \quad \text{(Structure: hierarchy $\leftrightarrow$ flexibility)}
\end{align}
\end{definition}

\paragraph{Interpretation.}
A point $(s, \sigma) = (0.7, 0.3)$ represents:
strategy leaning toward dynamism (70\%) but structure leaning toward hierarchy (30\%).
This is a \emph{misfit}: innovative goals with rigid structure.

\subsection{The Payoff Function: Formalizing Complementarity}

The payoff function captures three forces:
\begin{enumerate}
  \item \textbf{Misfit penalty:} Deviation between strategy and structure is costly
  \item \textbf{Innovation synergy:} High $s$ and high $\sigma$ together create value
  \item \textbf{Stability synergy:} Low $s$ and low $\sigma$ together create value
\end{enumerate}

\begin{definition}[Payoff Function]
\begin{equation}
  \Pi(s, \sigma; \Psi) = \underbrace{-\gamma (s - \sigma)^2}_{\text{Misfit penalty}}
                        + \underbrace{\alpha(\Psi) \cdot s \cdot \sigma}_{\text{Innovation synergy}}
                        + \underbrace{\beta(\Psi) \cdot (1-s) \cdot (1-\sigma)}_{\text{Stability synergy}}
  \label{eq:payoff}
\end{equation}
where $\gamma > 0$ measures misalignment cost, $\alpha(\Psi) > 0$ the return to innovation configuration,
and $\beta(\Psi) > 0$ the return to stability configuration.
\end{definition}

\paragraph{Why This Functional Form?}
The misfit term $-\gamma(s - \sigma)^2$ is zero when $s = \sigma$ and increasingly negative as $|s - \sigma|$ grows.
The innovation term $\alpha \cdot s \cdot \sigma$ is highest at $(1, 1)$ and captures complementarity: high $s$ needs high $\sigma$ to pay off.
The stability term $\beta \cdot (1-s) \cdot (1-\sigma)$ is highest at $(0, 0)$ with similar logic.

\paragraph{Example: Evaluating Configurations.}
With $\gamma = 1$, $\alpha = 2$, $\beta = 2$:
\begin{center}
\renewcommand{\arraystretch}{1.2}
\begin{tabular}{llcl}
\hline
\textbf{Configuration} & $(s, \sigma)$ & $\Pi$ & \textbf{Interpretation} \\
\hline
Stability Peak (A) & $(0, 0)$ & $0 + 0 + 2 = 2$ & Coherent, stable \\
Innovation Peak (B) & $(1, 1)$ & $0 + 2 + 0 = 2$ & Coherent, innovative \\
Valley (M) & $(0.5, 0.5)$ & $0 + 0.5 + 0.5 = 1$ & Coherent but low synergy \\
Misfit 1 & $(1, 0)$ & $-1 + 0 + 0 = -1$ & Innovation + hierarchy \\
Misfit 2 & $(0, 1)$ & $-1 + 0 + 0 = -1$ & Stability + flexibility \\
\hline
\end{tabular}
\end{center}

\subsection{Equilibria and Their Properties}

\begin{proposition}[Equilibrium Configurations]
The payoff function (\ref{eq:payoff}) has two local maxima at $A = (0, 0)$ with $\Pi_A = \beta$
and $B = (1, 1)$ with $\Pi_B = \alpha$, one saddle point at $M = (0.5, 0.5)$,
and two local minima at $(0, 1)$ and $(1, 0)$ with $\Pi = -\gamma$.
\end{proposition}

\paragraph{Which Peak Is Higher?}
$\Pi_A = \beta$ vs. $\Pi_B = \alpha$ depends on context $\Psi$:
if $\alpha(\Psi) > \beta(\Psi)$, the innovation peak B is higher;
if $\beta(\Psi) > \alpha(\Psi)$, the stability peak A is higher.

\subsection{The Coherence Index K}

\begin{definition}[Coherence Index]
\begin{equation}
  K(s, \sigma) = 1 - \frac{\min\{d((s,\sigma), A), d((s,\sigma), B)\}}{d_{max}}
  \label{eq:kformal}
\end{equation}
where $d$ is Euclidean distance and $d_{max} = 1/\sqrt{2}$.
\end{definition}

$K = 1$ at peaks A and B (perfect coherence); $K \to 0$ at the point equidistant from both peaks.

\subsection{The Quality Index Q}

\begin{definition}[Quality Index]
$Q(C^*) = \Pi(C^*; \Psi)$ where $C^*$ is the equilibrium configuration.
\end{definition}

Realized welfare depends on both: $W = Q(C^*) \cdot g(K)$ where $g(1) = 1$, $g(0) = 0$.

\subsection{The 144-Component Structure}

The two-axis model captures essentials but lacks richness.
Real decisions involve multiple utility types, dimensions, and time horizons.

\begin{definition}[Utility Components]
The full structure has 144 components $U = \{u_{i,v,d,t}\}$ indexed by:
\begin{itemize}
  \item $i \in \{INU, KNU, IDN\}$: Utility category (Individual, Collective, Identity) --- 3
  \item $v \in \{G, P\}$: Valence (Gain or Pain) --- 2
  \item $d \in \{F, E, P, S, D, Eco\}$: FEPSDE dimension --- 6
  \item $t \in \{0, 1, 2, 3\}$: Time horizon (instant, short, medium, long) --- 4
\end{itemize}
Total: $3 \times 2 \times 6 \times 4 = 144$ components.
\end{definition}

\begin{definition}[Aggregate Welfare Function]
\begin{equation}
  W = \sum_{i \in \{INU, KNU, IDN\}} \sum_{d \in FEPSDE} \sum_{t=0}^{3} 
      \delta_d(\Psi)^t \left[ G_{i,d,t} - \lambda_d(\Psi) \cdot P_{i,d,t} \right]
  \label{eq:fullwelfare}
\end{equation}
with discount factor $\delta_d(\Psi) \in (0, 1)$ and loss aversion $\lambda_d(\Psi) > 1$.
\end{definition}

\paragraph{Key Features.}
Discounting weights future less than present.
Loss aversion ($\lambda > 1$) weights pains more than gains, following Prospect Theory.
Dimension-specific parameters allow health to be discounted differently than money.

\subsection{The Complementarity Matrix}

With 144 components, the full complementarity matrix is $C \in \mathbb{R}^{144 \times 144}$
with 20,736 entries (10,440 unique by symmetry).

The matrix has block structure reflecting interactions between utility categories:
\begin{equation}
  C = \begin{pmatrix}
    C_{INU,INU} & C_{INU,KNU} & C_{INU,IDN} \\
    C_{KNU,INU} & C_{KNU,KNU} & C_{KNU,IDN} \\
    C_{IDN,INU} & C_{IDN,KNU} & C_{IDN,IDN}
  \end{pmatrix}
\end{equation}

\subsection{Computational Tractability: Why This Is Now Feasible}

\paragraph{The Apparent Problem.}
A $144 \times 144$ matrix with 10,440 unique parameters seems intractable:
\begin{itemize}
  \item How do we estimate so many parameters?
  \item How do we compute equilibria?
  \item How do we perform comparative statics?
\end{itemize}

This concern explains why comprehensive complementarity frameworks
were not developed earlier---the computational burden seemed prohibitive.

\paragraph{The Solution: Structure and Technology.}

\textbf{1. Exploiting Block Structure.}
The matrix is not arbitrary. It has hierarchical block structure:
\begin{itemize}
  \item Within-category blocks ($C_{INU,INU}$, etc.): Dense, strong interactions
  \item Cross-category blocks ($C_{INU,KNU}$, etc.): Sparser, weaker interactions
  \item Within-dimension blocks: Strongest complementarities
\end{itemize}

This structure reduces effective dimensionality dramatically.
A sparse representation with $\sim$500 non-negligible entries captures most variation.

\textbf{2. Low-Rank Approximation.}
Empirically, complementarity matrices are approximately low-rank:
\begin{equation}
  C \approx U \Sigma V^T
\end{equation}
where $U, V \in \mathbb{R}^{144 \times k}$ and $k \ll 144$.

With $k = 10-20$ factors, we capture 90\%+ of the variance.
This reduces estimation from 10,440 parameters to $\sim$3,000.

\textbf{3. Hierarchical Bayesian Estimation.}
Parameters are not estimated independently but hierarchically:
\begin{align}
  C_{ij} &\sim \mathcal{N}(\mu_{block(i,j)}, \sigma^2_{block}) \\
  \mu_{block} &\sim \mathcal{N}(\mu_{global}, \tau^2)
\end{align}

Information is pooled across similar entries, regularizing estimates
and enabling learning from limited data.

\textbf{4. Modern Computational Infrastructure.}
\begin{itemize}
  \item \textbf{GPU acceleration:} Matrix operations parallelize efficiently
  \item \textbf{Automatic differentiation:} Gradients computed without hand derivation
  \item \textbf{Probabilistic programming:} Bayesian inference automated (Stan, PyMC, NumPyro)
  \item \textbf{Cloud computing:} Massive parallel estimation feasible
\end{itemize}

\textbf{5. LLM-Assisted Implementation.}
Researchers can specify models in natural language:
\begin{quote}
``Estimate a hierarchical complementarity model with block structure
for INU/KNU/IDN categories, sparse cross-category effects,
and context-dependent parameters varying by institutional quality.''
\end{quote}

The code is generated automatically, errors are debugged interactively,
and results are explained in accessible language.

\paragraph{Concrete Example: Computing $K$.}
Suppose we have estimated $\hat{C}$ from data and derived $C^*(\Psi)$ from theory.
Computing the coherence index:
\begin{equation}
  K = 1 - \frac{\|\hat{C} - C^*\|_F}{\|C^*\|_F}
\end{equation}

\textbf{Before (2010):}
\begin{itemize}
  \item Derive $C^*$ analytically (intractable for complex specifications)
  \item Implement Frobenius norm calculation (error-prone)
  \item Debug numerical issues (weeks of work)
  \item Document and validate (more weeks)
\end{itemize}

\textbf{Now (2026):}
\begin{verbatim}
K = 1 - np.linalg.norm(C_hat - C_star, 'fro') / np.linalg.norm(C_star, 'fro')
\end{verbatim}

Or simply: ``Calculate the coherence index for this estimated complementarity matrix.''

\paragraph{The Transformation.}
What would have been a multi-year computational project in 2010
is now an afternoon's work in 2026.
This is not a minor convenience---it is what makes the framework \emph{usable}.

The historical evolution from manual matrix computation to AI-assisted analysis
is documented in Appendix~\ref{app:computationalhistory}, tracing the 60-year
journey that made this framework operationally feasible.

\subsection{From Mathematical Structure to Empirical Prediction}

\paragraph{What the Mathematics Provides.}
The formal apparatus above defines:
\begin{itemize}
    \item The \emph{structure} of payoff functions: $\Pi(s, \sigma; \Psi)$ with synergy and misfit terms
    \item The \emph{form} of welfare aggregation: 144 components with discounting and loss aversion
    \item The \emph{logic} of equilibrium: Fixed-point conditions for $C^*$
    \item The \emph{measures} of fit: $K$ for coherence, $Q$ for quality
\end{itemize}

\paragraph{What the Mathematics Cannot Provide.}
The mathematics is silent on parameter \emph{values}:
\begin{itemize}
    \item What is $\gamma$ (misfit penalty) for manufacturing versus service firms?
    \item How does $\alpha(\Psi)$ (innovation return) vary with institutional quality?
    \item What are the 10,440 entries of the complementarity matrix $C$?
    \item How does loss aversion $\lambda_d$ differ across FEPSDE dimensions?
\end{itemize}

These are \emph{empirical} questions that require calibration on data.

\paragraph{The LLM Limitation in Mathematical Contexts.}
Large language models can manipulate the mathematical formalism fluently.
Given the payoff function (\ref{eq:payoff}), an LLM can:
\begin{itemize}
    \item Derive first-order conditions for equilibrium
    \item Compute comparative statics symbolically
    \item Explain the economic interpretation of each term
    \item Generate code to solve the model numerically
\end{itemize}

But it cannot determine parameter values because:
\begin{enumerate}
    \item \textbf{Parameters are not in formulas.} The equation $\Pi = -\gamma(s-\sigma)^2 + \alpha s\sigma + \beta(1-s)(1-\sigma)$ contains $\gamma, \alpha, \beta$ as symbols. Their values exist in datasets, not texts.
    
    \item \textbf{Context-dependence is underspecified.} We write $\alpha(\Psi)$ to indicate dependence on context. But $\alpha(high\ trust) = ?$ and $\alpha(low\ trust) = ?$ require estimation.
    
    \item \textbf{Heterogeneity is averaged away.} The literature reports ``average'' effects. But prediction requires the full conditional distribution $\alpha | \Psi$.
\end{enumerate}

\paragraph{Example: Estimating Misfit Penalty $\gamma$.}
How costly is strategy-structure misalignment? The answer varies:
\begin{center}
\small
\begin{tabular}{lcc}
\toprule
\textbf{Context} & \textbf{Textual Evidence} & \textbf{Calibrated $\gamma$} \\
\midrule
Tech startups & ``Misfit is fatal'' & $\gamma \approx 2.5$ \\
Regulated utilities & ``Misfit is tolerated'' & $\gamma \approx 0.8$ \\
Family businesses & ``Relationships buffer misfit'' & $\gamma \approx 1.2$ \\
\bottomrule
\end{tabular}
\end{center}

The qualitative descriptions in the literature do not determine the quantitative parameters.
The numbers come from structural estimation on firm performance data---
the same calibration-on-data approach that Section~4.X develops for behavioral prediction.

\paragraph{The Calibration Program for Mathematical Economics.}
Just as behavioral economics moved from documenting ``effects exist'' to estimating
``effect sizes,'' mathematical economics must move from deriving ``equilibrium conditions''
to estimating ``equilibrium parameters.'' The 144-component structure above
is a \emph{schema} for organizing calibration, not a substitute for it.

The mathematical formalization provides:
\begin{enumerate}
    \item \textbf{Identification:} Which parameters can be estimated from which data?
    \item \textbf{Regularization:} How to pool information across similar contexts?
    \item \textbf{Prediction:} How to extrapolate to new contexts?
    \item \textbf{Validation:} How to test out-of-sample?
\end{enumerate}

Without calibration, the mathematics remains a language for describing possibilities.
With calibration, it becomes a tool for predicting actualities.

\subsection{Summary: The Mathematical Architecture}

\begin{center}
\renewcommand{\arraystretch}{1.3}
\begin{tabular}{lll}
\hline
\textbf{Object} & \textbf{Definition} & \textbf{Meaning} \\
\hline
$(s, \sigma)$ & Configuration in $[0,1]^2$ & What we want + how we organize \\
$\Pi(s, \sigma; \Psi)$ & Payoff function & Value of configuration in context \\
$C^*$ & Fixed point of $\Phi$ & Self-consistent structure \\
$K$ & Distance to nearest peak & Coherence/stability \\
$Q$ & Payoff at equilibrium & Quality of configuration \\
$W$ & 144-component aggregation & Full welfare \\
$\Psi$ & Context variable & Technology, institutions, norms \\
\hline
\end{tabular}
\end{center}

\subsection{What Comes Next}
\label{sec:ch8-next}

With the mathematical foundations established, subsequent chapters apply and extend the framework:

\begin{enumerate}[nosep]
    \item \textbf{Chapter 9:} Context as endogenous---how $\Psi$ modulates all parameters (CORE-WHEN)
    \item \textbf{Chapter 10:} Die Nutzenarchitektur---the full 144-component welfare structure (CORE-WHAT)
    \item \textbf{Chapter 11:} Awareness---how attention filters potential into effective utility (CORE-AWARE)
\end{enumerate}

%%%%%%%%%%%%%%%%%%%%%%%%%%%%%%%%%%%%%%%%%%%%%%%%%%%%%%%%%%%%%%%%%%%%%%%%%%%%%%%
% READING PATH BOX
%%%%%%%%%%%%%%%%%%%%%%%%%%%%%%%%%%%%%%%%%%%%%%%%%%%%%%%%%%%%%%%%%%%%%%%%%%%%%%%
\begin{tcolorbox}[
    colback=green!5!white,
    colframe=green!75!black,
    title={\textbf{Reading Path}},
    fonttitle=\bfseries
]
\begin{tabular}{@{}ll@{}}
\textbf{Previous:} & Chapter 7 \\
\textbf{Next:} & Chapter 9 \\
\textbf{Deep Dive:} & See Appendix G (Glossary) for definitions \\
\end{tabular}

\medskip
\textit{For readers short on time: Focus on the Intuition Box and Summary section.}
\end{tcolorbox}

